% WHICH HYPHEN OF A LIGATURE EARNS A DISCRETIONARY.
%
% A hyphen read in a paragraph is followed by an empty discretionary, so
% the line may break after it. Where the hyphen went into a LIGATURE, it is
% the LAST character taken into that ligature that decides -- the reference
% asks what the last character it put down was, not whether any of them was
% a hyphen. trip.tfm ligatures `-' with character 0 into `-', and
%
%   -\char0 -    ..\kern-1.00002
%                ..\rip - (ligature -^^@)
%                ..\rip -
%                ..\discretionary
%
% -- nothing after the ligature, whose last character was character 0; one
% after the plain `-'. `--' in cmr10 draws one because the SECOND of its
% two hyphens is a hyphen, which is the same rule.
%
% This engine gave the discretionary to a ligature any of whose characters
% was a hyphen. trip line 248 writes \--\--\char-0-A\-.
\catcode`\{=1 \catcode`\}=2
\tracingonline=1 \tracingparagraphs=1 \showboxdepth=99 \showboxbreadth=99
\font\rip=trip \rip \hyphenchar\rip=`-
\hyphenpenalty 88 \exhyphenpenalty 89
\hsize=1000pt \parindent=0pt \pretolerance=-1 \tolerance=10000 \hbadness=10000
\vsize=20000pt \output={\shipout\box255}
\message{[A a dash, char 0, a dash]}
\setbox1=\vbox{\noindent -\char0 -\par}
\showbox1
\message{[B a discretionary hyphen then a dash]}
\setbox1=\vbox{\noindent \--\par}
\showbox1
\message{[C two dashes]}
\setbox1=\vbox{\noindent --\par}
\showbox1
\message{[done]}
\end
